\section{来源审计：Scaling 不是一条曲线，而是一套 Operating System}

本讲由 Anthropic 联合创始人 Ben Mann 主讲，时间是 2025 年 2 月。课堂从 GPT-2/GPT-3 与 scaling laws，讲到 frontier training 的 distributed failure、loss spike、RLHF/Constitutional AI、安全评估、Responsible Scaling Policy（RSP）、mechanistic interpretability，以及 chat 产品与 API 的不同发布纪律。仓库保留官方封面和 1025 条时间戳字幕；本文用论文与 Anthropic 官方材料校验机制，不把演讲后的模型、收入、RSP 或产品状态倒写回课堂。

\begin{warningbox}{课堂口径与版本边界}
“过去一年十倍增长”“coding segment 三个月十倍”等数字只作为讲者当时的 workload 描述，不当作经审计财务结论。课堂使用 2025-era AI Safety Level（ASL）叙事；截至 2026 年 8 月 11 日，Anthropic 已在 2026 年 2 月 24 日发布 RSP v3.0，采用 capability thresholds、required safeguards 等更新框架。本文先忠实解释历史课程，再单独标记后续演化。
\end{warningbox}

\begin{importantbox}{Frontier model program 的四个闭环}
\begin{enumerate}
  \item \textbf{Scaling loop}：hypothesis $\rightarrow$ pilot runs $\rightarrow$ fit/forecast $\rightarrow$ compute allocation；
  \item \textbf{Training loop}：telemetry $\rightarrow$ anomaly $\rightarrow$ checkpoint/replay $\rightarrow$ fix/resume；
  \item \textbf{Safety loop}：pre-training $\rightarrow$ post-training $\rightarrow$ evaluation $\rightarrow$ safeguards/deployment；
  \item \textbf{Product loop}：chat experiment $\rightarrow$ behavior evidence $\rightarrow$ API contract $\rightarrow$ migration/deprecation。
\end{enumerate}
\end{importantbox}

\begin{knowledgebox}{术语消化：四种“规模”必须分开}
\textbf{Training scale} 是参数、token 与训练 FLOPs；\textbf{serving scale} 是请求量、context、latency 和 accelerator fleet；\textbf{capability scale} 是任务成功、泛化与 tool use；\textbf{risk scale} 是模型能力、访问方式和 safeguard failure 共同产生的潜在影响。收入或流量增长不能直接证明 capability，benchmark 增长也不能自动证明 safety。
\end{knowledgebox}

\subsection{增长为什么会改变工程假设}

快速需求增长不只意味着“多买 GPU”。Training 与 inference 争夺 capacity，API 用户要求稳定 contract，chat 用户期待更快迭代，安全团队要处理更多 adversarial traffic，on-call 则面对更高故障成本。Program 需要把 capability、compute、reliability 和 safeguards 放在同一资源模型中，而不是分别优化。

\lecturefigure{01-scaling-program.png}{Frontier scaling 是 demand、capability、compute 与 reliability 耦合的 program。}{本地字幕 00:00--01:00；概念重绘。}

\begin{knowledgebox}{读图：增长带来正反馈，也带来风险复利}
Demand 提供收入和真实任务；capability 提高 adoption；compute 支持训练与 serving；reliability 维持用户信任。任何一环失衡都会反向放大：capability release 没有 serving capacity 会造成排队，traffic 没有 abuse control 会扩大风险，训练挤占 inference 会伤害现有客户，API regression 则会把模型升级变成业务 outage。
\end{knowledgebox}

\subsection{本章小结}

Scaling program 的对象是完整生命周期。后文不会把“更大模型”当成单一答案，而会追问预测、训练、恢复、对齐、评估、发布与兼容性如何共同成立。

